\chapter{Anàlisi i requisits}
\label{cap:requisitos}

\section{Actors i context}

L'actor principal és un \textbf{estudiant autenticat}. El sistema no exposa dades entre usuaris: totes les files de negoci porten \texttt{user\_id} i polítiques RLS a Postgres. Un segon actor tècnic és el \textbf{worker d'ingestió}, que reclama treballs en cua (\texttt{claim\_next\_job}) dins del procés FastAPI. No hi ha rol d'administrador a l'MVP: qui desplega el projecte Supabase té accés de plataforma, fora de la UI.

\section{Cas d'ús principal}

El flux mínim implementat al producte és:

\begin{enumerate}
  \item Registre / login (correu i contrasenya; restabliment via enllaç). El proveïdor d'Auth admet OAuth; aquest MVP no l'exposa a la UI (annex~\ref{anexo:despliegue}).
  \item Crear una assignatura (onboarding si no n'hi ha cap).
  \item Pujar un PDF, seguir l'estat (\emph{en cua} / \emph{processant} / \emph{llest} / error) i, si cal, reindexar.
  \item Preguntar al xat (amb àmbit de documents opcional), llegir la resposta en streaming i obrir cites (nom de PDF + fragment).
  \item Generar flashcards o un qüestionari; conservar-los a la biblioteca; practicar o editar; importar/exportar JSON.
  \item Al planner, puntuar qualitat SM-2 de les targetes vençudes.
  \item Ajustar el perfil (nom visible, flag de RAG agentic).
\end{enumerate}

\begin{figure}[htbp]
\centering
\begin{tikzpicture}[
  font=\small,
  actor/.style={draw, rounded corners=2pt, align=center, inner sep=4pt, minimum width=2.2cm},
  uc/.style={draw, ellipse, align=center, inner sep=3pt, minimum width=3.3cm, minimum height=0.85cm}
]
\node[actor] (est) {Estudiant\\autenticat};
\node[uc, right=1.6cm of est] (asig) {Gestionar assignatures};
\node[uc, below=0.28cm of asig] (pdf) {Pujar i indexar PDF};
\node[uc, below=0.28cm of pdf] (xat) {Preguntar amb cites};
\node[uc, below=0.28cm of xat] (gen) {Generar i editar pràctica};
\node[uc, below=0.28cm of gen] (srs) {Repassar (SM-2)};
\node[actor, right=1.5cm of gen] (wrk) {Worker\\d'ingestió};
\draw (est.east) -- (asig.west);
\draw (est.east) -- (pdf.west);
\draw (est.east) -- (xat.west);
\draw (est.east) -- (gen.west);
\draw (est.east) -- (srs.west);
\draw (wrk.west) -- (pdf.east);
\end{tikzpicture}

\caption{Casos d'ús de l'MVP. L'estudiant inicia el cicle; el worker completa la indexació.}
\label{fig:usecases}
\end{figure}

\paragraph{Pre- i postcondició (cas d'ús de xat).}
Pre: hi ha almenys un document en estat llest a l'assignatura. Flux: l'estudiant envia una pregunta; el sistema classifica, recupera, genera amb cites i (si cal) desa el fil. Post: la UI mostra la resposta i les cites; no s'ha cridat l'LLM amb la clau al navegador.

\paragraph{Escenaris d'error.}
PDF sense text extraïble: l'etapa \texttt{parse} falla i l'estat del document és error (no s'inventa OCR). Job d'ingestió fallit: l'usuari pot reindexar (RF-04). JWT caducat o invàlid: l'API respon 401. Assignatura o document d'un altre usuari: 404 (no es filtra l'existència aliena). Context insuficient al xat: el prompt obliga a dir-ho, no a inventar.

\section{Requisits funcionals}

La taula~\ref{tab:rf} recull els requisits coberts per l'API i la UI actuals (extracte; el contracte complet és a l'annex~\ref{anexo:esquema}). La prioritat és MoSCoW: \emph{Must} per a l'MVP del cicle, \emph{Should} per a preferències, \emph{Could} per a depuració.

\begin{table}[htbp]
\centering
\caption{Requisits funcionals implementats i prioritat.}
\label{tab:rf}
\small
\begin{tabularx}{\textwidth}{lcX}
\toprule
Id & Prior. & Requisit \\
\midrule
RF-01 & Must & L'usuari gestiona assignatures (CRUD). \\
RF-02 & Must & L'usuari puja PDF associats a una assignatura (límit 25\,MB). \\
RF-03 & Must & El sistema indexa el PDF de manera asíncrona i n'informa l'estat. \\
RF-04 & Must & L'usuari pot reintentar una ingestió fallida (reindex). \\
RF-05 & Must & El xat respon fent servir només context recuperat, amb cites (fitxer, pàgines, snippet). \\
RF-06 & Must & El xat suporta streaming SSE i historial de fils per assignatura. \\
RF-07 & Must & Es generen flashcards i qüestionaris a partir de l'índex (lecture-focused). \\
RF-08 & Must & La biblioteca persisteix conjunts amb títol, origen, recompte i data d'actualització. \\
RF-09 & Must & Es poden crear, editar, esborrar, importar i exportar cartes i preguntes. \\
RF-10 & Must & El planner llista targetes vençudes i registra ressenyes SM-2. \\
RF-11 & Should & El perfil persisteix preferències de retrieval (p.~ex. agentic). \\
RF-12 & Could & Un mode debug (flag) permet executar rangs de pipeline. \\
\bottomrule
\end{tabularx}
\end{table}

\section{Requisits no funcionals}

\begin{itemize}
  \item \textbf{Seguretat.} Autenticació JWT Supabase (HS256); autorització per \texttt{user\_id} a cada cas d'ús; RLS amb \texttt{(select auth.uid())} al client; el navegador no posseeix claus d'LLM. No hi ha SLA ni pentest en aquest TFM.
  \item \textbf{Mida i recuperació.} Pujada màxima 25\,MB (\texttt{MAX\_UPLOAD\_SIZE\_BYTES}); cerca densa i lèxica amb $k=20$ i context final $k=8$. La latència d'ingestió i de xat \emph{no} es garanteix: depèn d'APIs externes.
  \item \textbf{Mantenibilitat.} Domini sense imports de FastAPI/Supabase/OpenAI; ports ABC; tests amb fakes; factory OpenAI/Gemini.
  \item \textbf{Observabilitat.} Estats de job i \texttt{pipeline\_stage\_runs}; telemetria OpenTelemetry opcional.
  \item \textbf{Qualitat de lliurament.} \texttt{ruff}, \texttt{mypy --strict}, pytest unitari, \texttt{tsc}, Vitest, Playwright smoke, Docker Compose, CI a GitHub Actions.
  \item \textbf{Usabilitat.} Una tasca per pantalla; interfície en castellà de tu; cites amb nom de fitxer; el hover no és l'única via de navegació.
\end{itemize}

\section{Restriccions}

\begin{itemize}
  \item PDF amb text extraïble (PyPDF); no hi ha OCR a l'MVP.
  \item Dependència d'APIs d'embeddings/LLM (cost i latència externs).
  \item Supabase allotjat (l'esquema viu a les migracions SQL del repositori).
  \item El flag \texttt{ENABLE\_RATE\_LIMIT} existeix (defecte \texttt{true} a la configuració), però \textbf{no hi ha middleware de quota implementat}. És deute conegut: el flag no s'ha d'interpretar com a protecció real.
\end{itemize}

\section{Traçabilitat}

La taula~\ref{tab:trazas} lliga objectius, requisits i on es verifica. No substitueix l'avaluació de retrieval (O6 / E3).

\begin{table}[htbp]
\centering
\caption{Traçabilitat objectiu -- requisit -- verificació.}
\label{tab:trazas}
\small
\begin{tabularx}{\textwidth}{llX}
\toprule
Obj. & RF & Verificació \\
\midrule
O1 & RF-01--12 & Aquest capítol; annex~\ref{anexo:esquema}. \\
O2 & RF-02--04 & Pipeline + worker (cap.~\ref{cap:desarrollo}); tests d'ingestió amb fakes. \\
O3 & RF-05, RF-06 & Política de retrieval i xat; tests de política i API. \\
O4 & RF-07--10 & Generació, biblioteca, SM-2 al domini; tests d'artefactes i SRS. \\
O5 & RF-01, RF-11 & JWT + UI; tests d'API amb fakes. \\
O6 & --- & Parcial: tests (taula~\ref{tab:calidad}), latència (taula~\ref{tab:bench}), cicle UI a captures (secció~\ref{sec:cicle-ui}). No tanca fidelitat. \\
\bottomrule
\end{tabularx}
\end{table}

\section{Històries d'acceptació}

O1--O5 es consideren coberts si existeix (i)~ruta API, (ii)~UI o CLI, (iii)~test unitari o d'integració. Els tests unitaris cobreixen SRS, política de retrieval, artefactes d'estudi, API amb fakes i agregació de benchmark. O6 no usa aquest criteri: el capítol~\ref{cap:resultados} declara què es mesura (programari, latència, recorregut qualitatiu) i què resta obert.
